\chapter{بيان ما نزل من الكتاب عاماً وأقسامه}

\section{مقدمة الدرس وتعريف العام والخاص}
الحمد لله رب العالمين، والصلاة والسلام على نبينا محمد وعلى آله وصحبه أجمعين، أما بعد؛ فهذا الدرس يتابع تقريرات الإمام الشافعي رحمه الله في كتاب «الرسالة» حول لغة العرب التي نزل بها القرآن، وأهمية فهم دلالات الألفاظ فيها كالعام والخاص والمطلق والمقيد.

وقد عرّف الإمام الشوكاني «العام» بأنه: اللفظ المستغرق لجميع ما وضع له بحسب وضع واحد دفعة (مثل: يا أيها الناس). أما «التخصيص» فهو: إخراج بعض ما كان داخلاً تحت العموم. وعرفه ابن الحاجب بأنه: قصر العام على بعض مسمياته.

\section{أقسام العام في القرآن الكريم}
قسّم الإمام الشافعي ما نزل من القرآن الكريم من حيث العموم والخصوص إلى عدة أقسام:

\subsection{القسم الأول: العام الذي يراد به العموم ولا يدخله التخصيص}
وهو اللفظ العام الذي يستغرق جميع أفراده ولا يخرج منه شيء. ومثّل له الشافعي بقوله تعالى: ﴿اللَّهُ خَالِقُ كُلِّ شَيْءٍ وَهُوَ عَلَىٰ كُلِّ شَيْءٍ وَكِيلٌ﴾، وقوله: ﴿خَلَقَ السَّمَاوَاتِ وَالْأَرْضَ﴾، وقوله: ﴿وَمَا مِن دَابَّةٍ فِي الْأَرْضِ إِلَّا عَلَى اللَّهِ رِزْقُهَا﴾. فكل شيء في السماوات والأرض خلقه الله، وكل دابة فرزقها على الله، فهذا عام باقٍ على عمومه ولا مخصص له أبداً.

\subsection{القسم الثاني: العام الذي يجمع بين العموم والخصوص}
وهو أن تأتي الآية بلفظ عام، لكن يراد ببعض أجزائها العموم وببعضها الآخر الخصوص. ومثّل لذلك بقوله تعالى: ﴿مَا كَانَ لِأَهْلِ الْمَدِينَةِ وَمَنْ حَوْلَهُم مِّنَ الْأَعْرَابِ أَن يَتَخَلَّفُوا عَن رَّسُولِ اللَّهِ وَلَا يَرْغَبُوا بِأَنفُسِهِمْ عَن نَّفْسِهِ﴾. 
فقوله: ﴿أَن يَتَخَلَّفُوا﴾ عام يراد به الخصوص، إذ لا يشمل العجزة والنساء والأطفال والمرضى، بل يخص من يطيق الجهاد. أما قوله: ﴿وَلَا يَرْغَبُوا بِأَنفُسِهِمْ عَن نَّفْسِهِ﴾ فهو عام يراد به العموم؛ فلا يجوز لأي مؤمن (سواء كان صحيحاً أو مريضا، رجلاً أو امرأة) أن يرغب بنفسه عن نفس رسول الله صلى الله عليه وسلم.

ومن ذلك أيضاً قوله تعالى: ﴿إِنَّا خَلَقْنَاكُم مِّن ذَكَرٍ وَأُنثَىٰ وَجَعَلْنَاكُمْ شُعُوبًا وَقَبَائِلَ لِتَعَارَفُوا﴾ فهذا عام يشمل كل البشر، ثم قال: ﴿إِنَّ أَكْرَمَكُمْ عِندَ اللَّهِ أَتْقَاكُمْ﴾ وهذا عام يراد به الخصوص؛ لأن التقوى لا تكون إلا لمن عقلها وكان مكلفاً بالغاً، ويخرج منها غير المكلفين كالصبي والمجنون، لقوله صلى الله عليه وسلم: «رفع القلم عن ثلاثة: عن النائم حتى يستيقظ، وعن الصبي حتى يبلغ، وعن المجنون حتى يفيق».

\subsection{القسم الثالث: العام الظاهر الذي يراد به الخصوص}
وهو لفظ مخرجه عام على كل الناس، ولكن يُراد به فئة معينة يُستدل عليها بالسياق أو بأدلة أخرى، ومن أمثلته:
\begin{itemize}
    \item قوله تعالى: ﴿الَّذِينَ قَالَ لَهُمُ النَّاسُ إِنَّ النَّاسَ قَدْ جَمَعُوا لَكُمْ فَاخْشَوْهُمْ﴾. لفظ (الناس) عام، لكنه هنا أُطلق وأُريد به الخصوص في المواضع الثلاثة: فالذين قالوا هم بضعة أشخاص (نفر قليل)، والمقول لهم هم النبي صلى الله عليه وسلم وصحابته، والذين جمعوا هم جيش أبي سفيان المنصرف من أحد. وبقية أهل الأرض لم يقولوا ولم يُقَل لهم ولم يجمعوا.
    \item قوله تعالى: ﴿يَا أَيُّهَا النَّاسُ ضُرِبَ مَثَلٌ فَاسْتَمِعُوا لَهُ ۚ إِنَّ الَّذِينَ تَدْعُونَ مِن دُونِ اللَّهِ لَن يَخْلُقُوا ذُبَابًا﴾. الخطاب هنا بلفظ (الناس) عام، ولكنه مخصوص بعبدة الأصنام، ويخرج منه المؤمنون الموحدون.
    \item قوله تعالى: ﴿ثُمَّ أَفِيضُوا مِنْ حَيْثُ أَفَاضَ النَّاسُ﴾. المراد بالناس هنا الحجاج الحاضرون في ذلك الموقف، وليس كل الناس على وجه الأرض.
    \item قوله تعالى: ﴿وَقُودُهَا النَّاسُ وَالْحِجَارَةُ﴾. استثنى القرآن الكريم من هذا العموم أهل الإيمان والأنبياء بقوله: ﴿إِنَّ الَّذِينَ سَبَقَتْ لَهُم مِّنَّا الْحُسْنَىٰ أُولَٰئِكَ عَنْهَا مُبْعَدُونَ﴾.
\end{itemize}

\section{أثر السياق والتركيب والإضافة في بيان المعنى}
دلالات الألفاظ في لغة العرب تختلف باختلاف ثلاثة أمور: السياق، والتركيب، والإضافة.

\textbf{أثر التركيب:} 
يغير المعنى تماماً مع بقاء نفس الألفاظ؛ فجملة «ما في البيت إلا محمد» تفيد الحصر المطلق بأن البيت خالٍ إلا منه، بينما جملة «ما محمد إلا في البيت» (وهي مكونة من نفس الألفاظ الخمسة) تفيد حصر مكانه هو، وقد يكون معه في البيت غيره.

\textbf{أثر الإضافة:} 
«يد النملة» تختلف عن «يد الفيل»، ولله المَثَل الأعلى، فإضافة الصفة للخالق سبحانه تقتضي الكمال والجلال الذي لا يشابهه فيه شيء من خلقه.

\textbf{أثر السياق:} 
وقد ضرب الشافعي له مثلاً بكلمة (القرية)، فالعرب تطلق القرية وتريد بها الأبنية، وتطلقها وتريد بها السُّكان، والسياق هو المُحَدِّد. قال تعالى: ﴿وَاسْأَلْهُمْ عَنِ الْقَرْيَةِ الَّتِي كَانَتْ حَاضِرَةَ الْبَحْرِ إِذْ يَعْدُونَ فِي السَّبْتِ﴾. السياق يبيّن أن المعتدي في السبت ليس الجدران والأبنية، بل سكان القرية. وكذلك قوله: ﴿رَبَّنَا أَخْرِجْنَا مِنْ هَٰذِهِ الْقَرْيَةِ الظَّالِمِ أَهْلُهَا﴾، فليس كل أهل القرية ظَلَمة لوجود مستضعفين مؤمنين فيها، لكن الأغلبية أو المتسلطين كانوا هم الظلمة. فهذه الوجوه تفهمها العرب بسليقتها، ومن جهل لغة العرب استشكل عليه فهم كتاب الله.